% Generated draft: requires mathematical and source review.
\documentclass[11pt]{article}
\usepackage[margin=25mm]{geometry}
\usepackage{fontspec}
\setmainfont{Latin Modern Roman}
\usepackage{amsmath,amssymb,mathtools,unicode-math}
\setmathfont{Latin Modern Math}
\usepackage{enumitem,needspace,xurl,hyperref}
\setlength{\parindent}{0pt}
\setlength{\parskip}{5pt}
\setlength{\emergencystretch}{3em}
\newcommand{\N}{\mathbb{N}}
\newcommand{\Z}{\mathbb{Z}}
\newcommand{\Q}{\mathbb{Q}}
\newcommand{\R}{\mathbb{R}}
\newcommand{\C}{\mathbb{C}}
\newcommand{\F}{\mathbb{F}}
\newcommand{\A}{\mathbb{A}}
\newcommand{\PP}{\mathbb{P}}
\newcommand{\E}{\mathbb{E}}
\newcommand{\eps}{\varepsilon}
\newcommand{\dd}{\,\mathrm{d}}
\begin{document}
\section*{UnsolvedMath ID: 30005429}
\textbf{Source problem code:} catalogue entry
\subsection{Definitions and mathematical statement}
\textbf{Recoverable scope and terminology.} The supplied record concerns regional modelling of Covid-19 transmission in Norway and Sweden during 2021. Its time-dependent data sources are temperature, virus-strain frequencies, Google behaviour metrics, and vaccination data. The stated method is computationally intensive Bayesian analysis of an epidemic model at regional level. No epidemic equations, probability model, prior distributions, observation model, parameter domains, or definition of transmission attribution are supplied.

\textbf{Research questions recoverable from the source.} Using those regional data, infer how behavioural patterns, virus strains, and temperature, interpreted as seasonality, affected transmission dynamics, and how much observed transmission can be attributed to the recorded behaviours. Vaccination data are included among the inputs, but their mathematical role is unspecified. This is an empirical inference objective; the supplied evidence does not determine a mathematical proposition or establish whether the intended effects are causal or associational.

\textbf{Reported empirical context.} The excerpt reports that the four Google metrics ``Work'', ``Transit'', ``Grocery and pharmacies'', and ``Retail and recreation'' capture \(40\%\text{--}70\%\) of all transmission. It reports that ``Work'' and ``Transit'' have greater impact on transmission than ``Grocery and pharmacies'' and ``Retail and recreation''. The meaning of ``capture'', the denominator and aggregation underlying this percentage range, and the precise comparison of impacts are not specified. These are reported findings, not mathematical assertions established by this draft. Remaining transmission is represented by nuisance parameters and may include transmission in households and various other settings. The ``Household'' metric was omitted because of a reported very strong negative correlation; the excerpt ends before identifying the other quantity in that correlation.

\textbf{Unspecified intervention problem.} The presentation heading also includes ``Optimal intervention strategies for minimizing total incidence during an epidemic''. No corresponding formulation is supplied: the epidemic dynamics, admissible interventions, intervention constraints or costs, initial and boundary conditions, time horizon, and definition of total incidence are absent. Consequently, no precise optimization problem can be recovered without introducing unsupported assumptions.

\textbf{Formulation limitation.} The supplied registry marks the statement as \texttt{unrecoverable} and explains that the heading is followed by empirical discussion rather than a standalone mathematical conjecture. Its separate \texttt{open} label is supplied curation, not verification that a particular mathematical claim is open. This statement-only draft records the research scope and unresolved qualifications; it does not supply a missing theorem or optimization model.
\subsection{Short English statement}
The source describes Bayesian analysis of regional Covid-19 transmission in Norway and Sweden during 2021 using behaviour, temperature, viral-strain, and vaccination data. It also names an intervention-optimization topic without giving its formulation. The excerpt is truncated, and neither topic supplies enough mathematical specifications for a precise theorem or optimization problem.
\subsection{Sources}
\begin{itemize}
\item Local UnsolvedMath registry: supplied complete source record, permanent ID \(30005429\), source problem code \texttt{catalogue entry}, title ``Behavioral and Seasonal Effects on COVID-19 Transmission'', Oberwolfach Reports Open Problems collection. The registry supplies the label \texttt{open}, the statement status \texttt{unrecoverable}, and the explanation that no separate clean mathematical statement is promoted.
\item Supplied report citation and original-report literature link: \emph{Design and Analysis of Infectious Disease Studies} (2023), Oberwolfach Reports, DOI 10.4171/owr/2023/8; page range not supplied by EMS. \url{https://doi.org/10.4171/owr/2023/8}. The excerpt attributes the presentation to Tom Britton, with joint work credited to Felix Gunther, Hilde Kjelgaard Brustad, Arnoldo Frigessi, and Lasse Leskela; these spellings follow the supplied excerpt. The separate \texttt{proposed\_by} field is empty.
\item Additional supplied literature source, labelled ``G\"unther--Brustad--Frigessi--Britton, 2026 publication record'': \url{https://su.diva-portal.org/smash/person.jsf?faces-redirect=true&language=en&pid=authority-person%3A84680&searchType=SIMPLE}. No specific publication title, theorem, or result from that record is supplied.
\item Literature assessment dated exactly \texttt{2026-08-22}: supplied dated curation, not independently verified. It labels the record open, describes it as combining a Covid-19 transmission-analysis project with an intervention topic rather than stating a single theorem, and says the report calls the analysis work in progress, so no fully resolved claim is identifiable. The supplied background requires independent verification of time-sensitive status and cited claims. No source visit or independent literature assessment is asserted here.
\end{itemize}
\end{document}
